\centerline{\Huge \bf  Домашнее задание }
\centerline{\Huge \bf Гомотетия}

\begin{flushright}
    \scriptsize{Правда думаете, что всё на этом кончится?}
\end{flushright}

\problem{1} Четырёхугольник разрезан диагоналями на четыре треугольника. Докажите, что точки пересечения медиан этих треугольников образуют параллелограмм.


\problem{2} На каждом из оснований $AD$ и $BC$ трапеции $ABCD$ построены вне трапеции равносторонние треугольники.Докажите, что отрезок, соединяющий третьи вершины этих треугольников, проходит через точку пересечения диагоналей трапеции.

\begin{wrapfigure}{r}{0.25\linewidth}
    \centering
    \includegraphics[width=\linewidth]{VII блок/Киви/lesson_4/image.png}
\end{wrapfigure}

\problem{3} На каждой из сторон треугольника $ABC$ построено по прямоугольнику так, что они попарно касаются вершинами (см. рисунок). Докажите, что прямые, соединяющие вершины треугольника $ABC$ с соответствующими вершинами треугольника $A_1B_1C_1$, пересекаются в одной точке.

\problem{4}\textit{Самые халявыные 3 балла.} Придумайте оригинальную задачу на гомотетию и решите её. Оригинальность условия будет проверяться!